\section{Předpisy a normy}

Hledání předpisů, norem a optimálních hodnot bylo použito dvou zdrojů - vyhledavač DuckDuckGo a digitální knihovna ASPI.

Pro nalezení doporučení optimálních hodnot veličin byl zadán dotaz \textit{"IEQ guidelines"} do vyhledavače DuckDuckGo. Výsledkem tohoto dotazu bylo několik odkazů, ze kterých tři odkazovaly na stejnou stránku, která byla databází různých směrnic a nařízení vydaných různými státy \parencite{infinityconferencegroupincHomeIEQGuidelines}. Tato stránka nese sice mnoho informací, ale data se stát od státu liší a v něktečých případech nejsou ani dostupná. Z tohoto důvodu byl dotaz změněn na \textit{"IAQ guidelines"}, který vyústil v zobrazení třech stránek (ve více odkazech), které se dle titulků zdály být relevantní. Z těchto stánek byl vybrán dokument WHO, který přesně definuje a předepisuje hladiny veličin týkajících se IAQ \parencite{worldhealthorganizationWHOGlobalAir2021}.

Pro hledání právních předpisů týkajících se vnitřního prostředí, byla použita digitální knihov-na ASPI. Pro hledání byl použit dotaz \textit{vnitřní prostředí} Výsledky byly filtrovány podle platných a účinných předpisů a seřazeny podle relevance.

V kategirii \textit{předpisy ČR} bylo nalezeno dohromady 145 výsledků. Z těchto výsledků bylo prvních 50 blíže analyzováno, jestli jsou v kontextu práce relevantní. Nakonec bylo vybráno 5 zdrojů.
 Prvním ze zdrojů je \textit{Metodický návod ZD25/2013 pro měření a hodnocení mikroklimatických podmínek na pracovišti a vnitřního prostředí staveb úplné a aktualní znění}, který vymezuje kritéria pro měření, způsoby měření \parencite{wolterskluwerMetodickyNavodZD252013}. 
Dále se nařízení, kterým je \textit{Nařízení 361/2007 Sb. , kterým se stanoví podmínky ochrany zdraví při práci úplné a aktualní znění}, kde lze najít přípustné hladiny různých veličin podle tříd práce \parencite{wolterskluwerNarizeni36120072007}.
Dále byla nalezena dvojice vyhlášek ministerstva zdravotnictví, \textit{Vyhláška 6/2003 Sb. , kterou se stanoví hygienické limity chemických, fyzikálních a biologických ukazatelů pro vnitřní prostředí pobytových místností některých staveb úplné a aktualní znění} a \textit{Vyhláška 304/2022 Sb. , kterou se mění vyhláška č. 6/2003 Sb., kterou se stanoví hygienické limity chemických, fyzikálních a biologických ukazatelů pro vnitřní prostředí pobytových místností některých staveb úplné a aktualní znění}. Tyto vyhlášky stanovují limitní požadavky na některé veličiny IEQ  \parencite{wolterskluwerVyhlaska620032003,wolterskluwerVyhlaska30420222022}.
V poslední řadě existuje \textit{Vyhláška 194/2007 Sb. , kterou se stanoví pravidla pro vytápění a dodávku teplé vody, měrné ukazatele spotřeby tepelné energie pro vytápění a pro přípravu teplé vody a požadavky na vybavení vnitřních tepelných zařízení budov přístroji regulujícími a registrujícími dodávku tepelné energie úplné a aktualní znění}, kde lze najít požadované stavy vnitřního prostředí pro různé typy budov a místností. Jedná se zde pouze o teplotu a vlhkost \parencite{wolterskluwerVyhlaska19420072007}.

